\chapter{The Spine}
\label{stone:spine}

When we first began growing our constructions, something curious happened. We noticed a pattern in the PhiBase coordinates. At the time, we ignored it. There was too much else to discover. Now it is time to go back.

Let us begin where we started. One short edge. Nothing else. Its PhiBase description is
\[
P(1,0).
\]
Now let it grow. Once,
\[
P(0,1).
\]
Again,
\[
P(1,1).
\]
Again,
\[
P(1,2).
\]
Again,
\[
P(2,3).
\]
Again,
\[
P(3,5).
\]
Again,
\[
P(5,8).
\]

Stop. Do not calculate anything. Simply look. The first coordinates are
\[
1,\;0,\;1,\;1,\;2,\;3,\;5.
\]
The second coordinates are
\[
0,\;1,\;1,\;2,\;3,\;5,\;8.
\]
Do those numbers seem familiar? They should. The Fibonacci sequence has been quietly walking beside us since the moment we learned how to grow a golden triangle. We simply hadn't noticed it yet.

\section*{The Geometry Already Knew}

This is one of the happiest moments in mathematics. We did not insert the Fibonacci numbers. We did not force them into the arithmetic. The geometry produced them naturally. Every time a short edge grew into a long edge, and every long edge grew into one short edge and one long edge, the Fibonacci sequence took one more step. The arithmetic was simply keeping score.

\section*{Two Directions}

Yesterday we discovered something else. The conjugate introduced PhiBase numbers like
\[
P(5,-3),
\]
and
\[
P(3,-2).
\]
At first they seemed completely new. Now look more carefully. Write them beside the positive sequence:
\[
P(1,0),\; P(0,1),\; P(1,1),\; P(1,2),\; P(2,3),\; P(3,5)
\]
and
\[
P(2,-1),\; P(3,-2),\; P(5,-3),\; P(8,-5).
\]
The same Fibonacci numbers appear again. Only the direction has changed.

The Spine does not stop at \(1\). It continues forever in both directions. Growth. And descent. Exactly like the number line itself.

\section*{A Different Way To Think About Powers}

Until now, powers of \(\varphi\) may have seemed like repeated multiplication. There is another way to see them. Each step along the Spine is simply another growth. Move upward one rung. Multiply by \(\varphi\). Move downward one rung. Divide by \(\varphi\). The arithmetic has become movement. The powers of \(\varphi\) have become places.

\section*{A Builder's View}

Imagine a ladder. Each rung is one perfectly exact construction. Climbing upward never changes the rule. Climbing downward never changes the rule. Every rung knows the one above it. Every rung knows the one below it.

That ladder is the Fibonacci Spine. It is not merely a list of numbers. It is the backbone of PhiBase.

\section*{Builder's Notebook}

Beginning with
\[
P(1,0),
\]
grow upward eight times. Now begin again. This time descend eight times. Write both sequences side by side. Do not look for formulas. Simply observe. What remains the same? What changes? Can you describe the Spine without mentioning Fibonacci at all?

\section*{Looking Ahead}

The Spine gives us a wonderfully organized collection of exact PhiBase numbers. Real constructions, however, are not confined to a single ladder. Between every pair of Spine entries lie thousands of other perfectly good PhiBase numbers. How should we organize them? That question leads us to the Dense Grid.
